\section{Which finite \'etale algebras occur?}

Let $K$ be a field of characteristic zero.  A polynomial map
\[
 F:\A^m_K\longrightarrow\A^m_K
\]
is a \emph{Keller map} if $\det DF\in K^*$.  Its geometric degree is
\[
 \gdeg(F)=[K(x_1,\ldots,x_m):K(F_1,\ldots,F_m)].
\]
A Keller map is \'etale and hence quasi-finite.  Every fiber over a $K$-point
is therefore a finite \'etale $K$-scheme, possibly empty.

\begin{definition}[Keller fiber]\label{def:keller-fiber}
A nonzero finite \'etale $K$-algebra $A$ is a \emph{Keller fiber} if there are
a Keller map $F:\A^m_K\to\A^m_K$ and a point $y\in\A^m(K)$ such that
\[
 F^{-1}(y)\simeq\Spec A,
 \qquad \dim_KA=\gdeg(F).
\]
The equality of the two degrees is the fullness condition.  It says that the
special fiber contains all generic inverse sheets.  Without it, a fiber may be
shorter because some sheets have escaped through the nonproperness boundary.
\end{definition}

The construction used in the proof can be read from left to right:
\[
 \boxed{
 \text{squarefree }P(T)
 \ \longmapsto\
 G(S)=P(a+S)-P(a)
 \ \longmapsto\
 \widetilde F_G^{-1}(y_{P,a})\simeq\Spec K[T]/(P).}
\]
The last arrow is an isomorphism of schemes.  It is not merely a matching of
$K$-points or a count after passing to an algebraic closure.

\begin{lemma}[Degree two is impossible, over every characteristic-zero field]
\label{lem:degree-two}
Let $K$ be a field of characteristic zero, let $m\geq1$, and let
\[
 F=(F_1,\ldots,F_m)\in K[X_1,\ldots,X_m]^m
 \quad\text{satisfy}\quad
 \det\left(\frac{\partial F_i}{\partial X_j}\right)\in K^*.
\]
Then the induced function-field extension cannot have degree two:
\[
 \left[
 K(X_1,\ldots,X_m):K(F_1,\ldots,F_m)
 \right]\ne2.
\]
\end{lemma}

\begin{proof}
We give the scalar-extension and descent steps because the external theorem
used here is most conveniently cited over $\mathbb C$.

\smallskip
\noindent\emph{The precise external input.}
Campbell's unnumbered theorem
\cite[Theorem, p.~244]{Campbell} states that a polynomial map
$H:\mathbb C^m\to\mathbb C^m$ has a polynomial inverse if its Jacobian
determinant never vanishes and the induced extension
\[
 \mathbb C(H_1,\ldots,H_m)\subset
 \mathbb C(X_1,\ldots,X_m)
\]
is normal.  In characteristic zero ``normal'' here is equivalent to
``Galois,'' because the finite function-field extension is separable.
Razar and, independently, Wright later gave algebraic treatments of this
Galois case \cite{Razar,Wright}; the proof below needs only Campbell's stated
complex theorem.

\smallskip
\noindent\emph{Generic degree survives scalar extension.}
Assume for contradiction that the displayed degree over $K$ is two.  Let
$K_0\subset K$ be the subfield generated over $\mathbb Q$ by the finitely
many coefficients of the $F_i$, and write $F_0$ for the same formulas over
$K_0$.  The nonzero Jacobian determinant makes
\begin{equation}
 \varphi_0:K_0[Y_1,\ldots,Y_m]\longrightarrow
 K_0[X_1,\ldots,X_m],
 \qquad Y_i\longmapsto F_i,
 \label{eq:degree-two-coordinate-map}
\end{equation}
injective and generically finite.  Indeed, the Jacobian criterion in
characteristic zero makes the $F_i$ algebraically independent, so both
fraction fields in \eqref{eq:degree-two-coordinate-map} have transcendence
degree $m$ and their extension is finite.  If its degree is $d_0$, there is a
nonempty principal open $U_0$ of the target on which $F_0$ is finite locally
free of rank $d_0$: first clear denominators in the finite extension of
fraction fields, and then apply generic freeness.  For every field extension
$E/K_0$, the base change over the nonempty open $(U_0)_E$ is again finite
locally free of rank $d_0$.  Taking generic fibers gives
\begin{equation}
 [E(X_1,\ldots,X_m):E(F_1,\ldots,F_m)]=d_0.
 \label{eq:degree-two-base-change}
\end{equation}
Applying \eqref{eq:degree-two-base-change} first to $E=K$ shows $d_0=2$.

\smallskip
\noindent\emph{Why degree two gives Campbell's field hypothesis.}
Because $K_0$ is finitely generated over $\mathbb Q$, choose an embedding
$K_0\hookrightarrow\mathbb C$.  Equation
\eqref{eq:degree-two-base-change} with $E=\mathbb C$ says that
\[
 \mathbb C(X_1,\ldots,X_m)/
 \mathbb C(F_1,\ldots,F_m)
\]
has degree two.  It is separable in characteristic zero.  Every separable
quadratic field extension is normal (the second root of the minimal
polynomial of an element is its trace minus that element), hence it is
Galois.  Thus the complexified map satisfies exactly the hypotheses of
Campbell's theorem and is a polynomial automorphism.

\smallskip
\noindent\emph{Faithfully flat descent of the inverse.}
The base change of \eqref{eq:degree-two-coordinate-map} along
$K_0\hookrightarrow\mathbb C$ is therefore an isomorphism.  Since
$\mathbb C$ is faithfully flat over $K_0$, tensoring the kernel and cokernel
of $\varphi_0$ with $\mathbb C$ detects their vanishing.  Hence $\varphi_0$
itself is an isomorphism.  After base change to $K$, the original $F$ is a
polynomial automorphism, so
$K(F_1,\ldots,F_m)=K(X_1,\ldots,X_m)$ and its function-field degree is one.
This contradicts the assumed degree two.
\end{proof}

\begin{theorem}[Prescribed finite-\'etale realization]
\label{thm:main}
Let $P\in K[T]$ be squarefree of degree $N\geq3$.  Choose $a\in K$ outside
the finite zero set of $P'P'''$, and put
\[
 G(S)=P(a+S)-P(a).
\]
Let $F_G=(\Pi,B,C)$ be the quadratic-gauge map of
\eqref{eq:gauge-map}, and define
\[
 \widetilde F_G=L\circ F_G,
 \qquad
 L(\Pi,B,C)=\left(\Pi,-\frac{B}{2},C\right).
\]
Then
\[
 \det D\widetilde F_G=1,
 \qquad \gdeg(\widetilde F_G)=N,
 \qquad \max_i\deg (\widetilde F_G)_i\leq 6N+2.
\]
The Galois closure of the induced degree-$N$ function-field extension has
geometric and arithmetic Galois group $S_N$.
At
\begin{equation}
 y_{P,a}=\left(1,0,-\frac{2P(a)}{P'(a)}\right)
 \label{eq:main-target}
\end{equation}
its full fiber is
\[
 \boxed{\widetilde F_G^{-1}(y_{P,a})
 \simeq\Spec K[T]/(P).}
\]
More explicitly, for every commutative $K$-algebra $R$ there is a bijection
\[
 \operatorname{Hom}_{K\text{-alg}}(K[T]/(P),R)
 \simeq
 \{u\in R^3:\widetilde F_G(u)=y_{P,a}\},
\]
natural in $R$.

Consequently, every finite \'etale $K$-algebra of rank at least three occurs
as a full fiber of a Jacobian-one map in affine three-space.
\end{theorem}

\begin{proof}
The polynomials $P'$ and $P'''$ are nonzero in characteristic zero when
$N\geq3$.  Their product has only finitely many roots, while $K$ is infinite,
so an admissible $a$ exists.  The linear and cubic Taylor coefficients of $G$
are $P'(a)$ and $P'''(a)/3!$, and the quadratic gauge of
Section~\ref{sec:gauge} is therefore defined.

At the target \eqref{eq:main-target}, its inverse polynomial is
\[
 G(S)-\frac{P'(a)}2\left(-\frac{2P(a)}{P'(a)}\right)
 =P(a+S).
\]
Translation preserves squarefreeness, so
Corollary~\ref{cor:squarefree-gauge}
identifies the full fiber with $\Spec K[S]/(P(a+S))$.

The translation back to the original polynomial is explicit:
\[
 K[S]/(P(a+S))\longrightarrow K[T]/(P(T)),
 \qquad \overline S\longmapsto\overline T-a,
\]
with inverse $\overline T\mapsto\overline S+a$.  Composing this with the
natural equivalence of Corollary~\ref{cor:squarefree-gauge} gives, for every
commutative $K$-algebra $R$,
\[
 \operatorname{Hom}_{K\text{-alg}}(K[T]/(P),R)
 \simeq \widetilde F_G^{-1}(y_{P,a})(R).
\]
Thus the isomorphism includes the complete scheme structure and commutes with
base change.

The linear map $L$ fixes the chosen target because its second coordinate is
zero.  It multiplies the Jacobian by $-1/2$ and changes neither geometric
degree nor coordinate degree.  Corollary~\ref{cor:effective-gauge} therefore
gives the determinant-one statement and the bound $6N+2$.  Since both the
quotient algebra and the generic inverse problem have degree $N$, the fiber is
full.  The symmetric-monodromy assertion follows from
Theorem~\ref{thm:symmetric-monodromy}.

For a general finite \'etale algebra of rank at least three, use the
monogenicity lemma in Section~\ref{sec:realization} and apply the polynomial
case just proved.
\end{proof}

\begin{theorem}[Classification of full Keller fibers]
\label{cor:rank-classification}
For $N\geq1$, let $\mathsf{FullFib}_{3,N}(K)$ be the replete full subcategory
of finite \'etale $K$-schemes whose objects are the schemes isomorphic to
full fibers
\[
 F^{-1}(y),\qquad
 F:\A^3_K\longrightarrow\A^3_K,\quad
 \gdeg(F)=N,\quad y\in\A^3(K).
\]
Then
\[
 \mathsf{FullFib}_{3,N}(K)=
 \begin{cases}
  \mathsf{FinEt}^{\,N}_K,&N=1\text{ or }N\geq3,\\
  \varnothing,&N=2,
 \end{cases}
\]
where $\mathsf{FinEt}^{\,N}_K$ is the full subcategory of rank-$N$ finite
\'etale $K$-schemes.  In particular, for $N\geq3$, a finite $K$-scheme occurs
as a full fiber of a geometric-degree-$N$ Keller map in $\A^3_K$ if and only
if it is finite \'etale of rank $N$.

Equivalently on coordinate algebras, every rank-$N$ finite \'etale
$K$-algebra occurs for $N=1$ or $N\geq3$, and none occurs for $N=2$.
Consequently, the possible nonzero ranks of full Keller fibers are exactly
\[
 \boxed{1,3,4,5,\ldots}.
\]
\end{theorem}

\begin{proof}
The Jacobian condition makes $F$ \'etale.  Hence every fiber over a
$K$-point is a finite \'etale $K$-scheme, and the fullness equality makes its
rank equal to $\gdeg(F)=N$.  This proves the necessary direction in every
degree.  The identity map realizes the unique rank-one object.
Theorem~\ref{thm:main} realizes every rank-$N$ finite \'etale object in
affine three-space for every $N\geq3$, while
Lemma~\ref{lem:degree-two} shows that no geometric-degree-two Keller map
exists.  This last exclusion concerns the generic function-field degree
before any target point is selected.  Thus an empty special fiber cannot
create a rank-two object; it is not a nonzero finite \'etale scheme at all.
Nor can a shorter special fiber of a map of some other geometric degree:
such a fiber fails the defining fullness equality and is not an object of
$\mathsf{FullFib}_{3,2}(K)$.  Empty and non-full fibers therefore have no
bearing on the asserted rank set.
\end{proof}

\begin{remark}[What the categorical statement remembers]
The category $\mathsf{FullFib}_{3,N}(K)$ is the essential image inside finite
\'etale $K$-schemes: it remembers the resulting fiber scheme, not a chosen
presentation by a map--target pair $(F,y)$.  Different pairs may realize the
same object.  If one works with finite \'etale algebras rather than schemes,
the corresponding statement is contravariant under $\Spec$.
\end{remark}
